\section{Desarrollo Histórico}

\begin{frame}{El Mercurio en la Antigüedad}
  \begin{block}{Un Elemento Ancestral}
    Conocido desde al menos 1500 a. C., se han encontrado viales en antiguas
    tumbas egipcias. Fue reconocido por prácticamente todas las grandes
    civilizaciones antiguas.
  \end{block}

  \begin{columns}[T]
    \begin{column}{0.5\textwidth}
      \begin{exampleblock}{Etimología: Azogue}
        Del latín \textit{argentum vivum}, que significa ``plata viva",
        reflejando su naturaleza metálica y móvil.
      \end{exampleblock}
    \end{column}
    \begin{column}{0.5\textwidth}
      \begin{alertblock}{Símbolo: \ce{Hg}}
        De su nombre griego, \textit{Hydrargyrum}, que significa ``plata de
        agua". Este fue latinizado posteriormente y adoptado para su símbolo
        químico.
      \end{alertblock}
    \end{column}
  \end{columns}
\end{frame}

\begin{frame}{El Corazón de la Alquimia}
  \begin{block}{La \textit{Prima Materia}}
    Durante siglos, los alquimistas creyeron que el Mercurio era la
    \textbf{Prima Materia} (Primera Materia) a partir de la cual se formaban
    todos los demás metales.
  \end{block}

  \begin{alertblock}{Los Tres Principios (\textit{Tria Prima})}
    El Mercurio era considerado una de las tres sustancias fundamentales del
    mundo material, junto con el Azufre y la Sal.
    \begin{itemize}
      \item \textbf{Mercurio}: El principio de fusibilidad y volatilidad (el espíritu).
      \item \textbf{Azufre}: El principio de inflamabilidad (el alma).
      \item \textbf{Sal}: El principio de incombustibilidad (el cuerpo).
    \end{itemize}
  \end{alertblock}
\end{frame}

\begin{frame}{La Teoría Alquímica de la Transmutación}
  \begin{block}{El Objetivo: Crear Oro}
    Los alquimistas teorizaban que, purificando y combinando Mercurio y Azufre
    en diferentes proporciones, podían transmutarlos en otros metales.
  \end{block}

  \begin{itemize}
    \item El oro era considerado el metal más perfecto y estable, requiriendo el
      equilibrio más perfecto de estos principios.
    \item La legendaria \textbf{Piedra Filosofal} era la sustancia que se creía
      que permitía esta transmutación perfecta.
    \item Esta búsqueda no era solo química, sino también profundamente
      filosófica y espiritual.
  \end{itemize}
\end{frame}

\begin{frame}{De la Alquimia a la Química}
  \begin{alertblock}{El Fin de una Era}
    El papel central del mercurio en las teorías de transmutación fue refutado
    por el auge de la química cuantitativa y experimental.
  \end{alertblock}

  \begin{block}{Antoine Lavoisier (1743-1794)}
    \begin{itemize}
      \item A través de experimentos rigurosos, incluyendo la calcinación
        (oxidación) y reducción del óxido de mercurio, Lavoisier estableció el
        concepto moderno de elemento.
      \item Demostró que el mercurio no era un principio, sino un verdadero
        elemento que se combinaba con el oxígeno del aire.
      \item Este trabajo fue fundamental para la Revolución Química y la
        desaparición de la alquimia como disciplina científica.
    \end{itemize}
  \end{block}
\end{frame}
